\section{Introduction}
\label{sec:introduction}

The initial stage of architectural design, known as massing, involves defining the overall three-dimensional shape and volume of a building. This phase is critical as it establishes the project's relationship with its site, its functional organization, and its fundamental economic viability. Traditionally, generating and evaluating massing options is a manual, time-consuming process that relies heavily on the designer's intuition and experience. This reliance creates a bottleneck in the design workflow and can limit the exploration of a wide range of contextually sensitive and programmatically optimized solutions.

While computational methods have evolved from algorithmic optimization to data-driven and generative AI approaches, existing techniques cannot simultaneously handle functional requirements, geometric precision, and visual context. The absence of suitable multi-modal datasets has hindered progress toward unified massing generation systems.

This paper presents three key contributions: (1) \textbf{CoMa-20K}, a novel dataset with 20,000 sites combining massing geometries, functional requirements, and urban context views; (2) a \textbf{VLM-based benchmark} formulating massing as a conditional generation task, evaluated through both fine-tuning and zero-shot approaches; and (3) a \textbf{comprehensive evaluation} revealing fundamental challenges in geometric reasoning and contextual adaptation.

Our work establishes a foundation for data-driven architectural design, demonstrating both the feasibility and complexity of automated massing generation.